\section{Motivation}
\label{sec:motivation}

Many NoSQL stores offer fast single-key operations but no transactions spanning several keys, let
alone several stores. Applications that need to update data atomically across such stores rely on
an external coordinator running two-phase commit (2PC), and on a \emph{shim} in front of each data
source that isolates the branches executing there. The natural choice for this shim is strict
two-phase locking (S2PL): a branch locks every key it accesses and releases its locks only once
it has committed.

Under 2PC, however, the lifetime of a branch does not end with its execution. Once it has sent its
last operation, it still has to be prepared, vote, wait for the votes of the other branches of its
global transaction, and receive the coordinator's decision before it can commit. Its writes are
already final when its execution ends: nothing after that point can change them. Yet S2PL keeps its
locks for the whole window, which spans the prepare round trip, the slowest branch of the global
transaction and the decision round trip. Every conflicting transaction waits through this window
for values that will no longer change, and under skewed access long queues form behind the hot
keys, so throughput is bounded by commit latency rather than by execution.

The shim can instead \emph{speculate}: let a waiting transaction read the final but uncommitted
values of the transactions ahead of it as soon as they have finished executing, rather than once
they have committed. Correctness is preserved by \emph{commitment ordering} (CO): a transaction may
only commit after every conflicting transaction ahead of it has resolved, so commits follow the
conflict order. Since the local history of every shim is CO, the global history is serializable,
without any global concurrency control beyond 2PC itself.

Speculation is not free. A transaction that speculated on another one depends on it: if the latter
aborts, the former must abort too, and so must every transaction that speculated on it in turn.
These \emph{cascading aborts} grow with the amount of speculation. The more transactions read on top
of uncommitted ones, the longer the chains of dependents and the more work a single abort destroys.
Beyond some contention level, the work lost to cascades exceeds the waiting saved, and speculation
falls behind the non-speculative shim. This raises the question this project addresses:
\emph{can a shim speculate only as much as it pays off, and so get the best of both worlds?}

\section{Contributions}
\label{sec:contributions}

\begin{itemize}[leftmargin=*]
  \item A speculative, CO-compliant shim for NoSQL stores that participates in 2PC as a resource
        manager, together with a non-speculative S2PL variant built from the same code, so that
        the two differ only in when a lock is granted.
  \item A cap on speculation, the number of transactions that may speculate on a single
        uncommitted writer, and an adaptive controller that tunes it at run time from the shim's
        commit rate, moving between full speculation and none.
  \item An evaluation of the three shims on a micro-benchmark, across key skew, transaction length,
        number of branches and thread counts.
\end{itemize}

\section{Outline}
\label{sec:outline}

\Cref{ch:design} presents the design of the speculative shim and of the adaptive cap.
\Cref{ch:implementation} describes their implementation. \Cref{ch:evaluation} compares the
speculative, non-speculative and adaptive shims, and \cref{ch:conclusion} concludes and discusses
future work.

\section{Related Work}
\label{sec:related-work}
